\documentclass[11pt,letterpaper]{article}

\usepackage[top=1in, bottom=1in, left=1in, right=1in, headheight=14pt]{geometry}
\usepackage{fontspec}
\setmainfont{DejaVu Serif}
\setsansfont{DejaVu Sans}
\setmonofont{DejaVu Sans Mono}

\usepackage{xcolor}
\usepackage{titlesec}
\usepackage{fancyhdr}
\usepackage{enumitem}
\usepackage{hyperref}
\usepackage{booktabs}
\usepackage{tabularx}
\usepackage{longtable}
\usepackage{colortbl}
\usepackage{multirow}
\usepackage{array}
\usepackage{parskip}
\usepackage{tcolorbox}
\usepackage{graphicx}
\usepackage{lastpage}

% === COLOR SCHEME: Technology / Hardware / Retrocomputing ===
\definecolor{primary}{HTML}{1A237E}       % Deep indigo — architecture
\definecolor{secondary}{HTML}{0277BD}     % Electric blue — implementation
\definecolor{tertiary}{HTML}{00695C}      % Teal — integration
\definecolor{accent}{HTML}{1565C0}        % Accent blue
\definecolor{lightprimary}{HTML}{E8EAF6}
\definecolor{lightsecondary}{HTML}{E1F5FE}
\definecolor{lighttertiary}{HTML}{E0F2F1}
\definecolor{rowgray}{HTML}{F5F5F5}
\definecolor{bordergray}{HTML}{BDBDBD}
\definecolor{subtlegray}{HTML}{757575}
\definecolor{darktext}{HTML}{212121}

% === SECTION FORMATTING ===
\titleformat{\section}
  {\Large\bfseries\sffamily\color{primary}}
  {\thesection}{1em}{}
  [\vspace{-0.5em}\textcolor{bordergray}{\rule{\textwidth}{0.4pt}}]

\titleformat{\subsection}
  {\large\bfseries\sffamily\color{secondary}}
  {\thesubsection}{1em}{}

\titleformat{\subsubsection}
  {\normalsize\bfseries\sffamily\color{tertiary}}
  {\thesubsubsection}{1em}{}

\titlespacing*{\section}{0pt}{1.5em}{0.8em}
\titlespacing*{\subsection}{0pt}{1.2em}{0.5em}
\titlespacing*{\subsubsection}{0pt}{0.8em}{0.3em}

% === CUSTOM ENVIRONMENTS ===
\newcommand{\stageheader}[2]{%
  \noindent\colorbox{#2}{\makebox[\textwidth][l]{%
    \textcolor{white}{\bfseries\sffamily\hspace{6pt}#1\hspace{6pt}}}}%
  \vspace{0.5em}\par%
}

\tcbuselibrary{breakable,skins}

\newtcolorbox{asciibox}{
  colback=rowgray, colframe=bordergray,
  arc=1pt, boxrule=0.5pt,
  left=6pt, right=6pt, top=4pt, bottom=4pt,
  fontupper=\small\ttfamily,
  breakable
}

\newtcolorbox{quotebox}{
  colback=lightprimary, colframe=primary,
  arc=2pt, boxrule=0.5pt,
  left=12pt, right=12pt, top=8pt, bottom=8pt,
  breakable
}

\newcommand{\metafield}[2]{\textbf{\sffamily #1} #2\par}

% === TABLE HELPERS ===
\newcolumntype{L}[1]{>{\raggedright\arraybackslash}p{#1}}
\newcolumntype{C}[1]{>{\centering\arraybackslash}p{#1}}
\newcommand{\tableheader}[1]{\textbf{\sffamily\textcolor{white}{#1}}}

% === HEADERS / FOOTERS ===
\pagestyle{fancy}
\fancyhf{}
\fancyhead[L]{\small\sffamily\textcolor{subtlegray}{GPU $\cdot$ GSD-OS $\cdot$ Observability}}
\fancyhead[R]{\small\sffamily\textcolor{subtlegray}{GSD Mission Package}}
\fancyfoot[C]{\small\sffamily\textcolor{subtlegray}{Page \thepage\ of \pageref{LastPage}}}
\renewcommand{\headrulewidth}{0.4pt}
\renewcommand{\footrulewidth}{0pt}

% === HYPERLINKS ===
\hypersetup{
  colorlinks=true,
  linkcolor=secondary,
  urlcolor=secondary,
  pdfauthor={GSD Ecosystem / Tibsfox},
  pdftitle={GPU Graphics, Compute, and GSD-OS Observability -- Mission Package},
  pdfsubject={Vision to Mission Pipeline},
}

% ================================================================
\begin{document}
% ================================================================

% ----------------------------------------------------------------
% TITLE PAGE
% ----------------------------------------------------------------
\thispagestyle{empty}
\begin{center}
\vspace*{2.5cm}

{\small\sffamily\textcolor{primary}{\textbf{GSD RESEARCH MISSION PACKAGE --- FLEET ACTIVATION}}}

\vspace{0.3cm}
\textcolor{bordergray}{\rule{8cm}{0.4pt}}

\vspace{1.5cm}
{\fontsize{26}{32}\selectfont\bfseries\sffamily\textcolor{primary}{GPU Graphics, Compute \& GSD-OS\\[0.3em]Observability Ecosystem}}

\vspace{0.8cm}
{\Large\sffamily\textcolor{secondary}{OpenGL/Vulkan $\cdot$ CUDA/OpenCL $\cdot$ SoC Audio $\cdot$ Amiga Expansion\\[0.2em]XScreenSaver Engineering $\cdot$ GSD-OS Integration $\cdot$ OTel/Prometheus/Grafana}}

\vspace{0.5cm}
{\large\sffamily\itshape\textcolor{tertiary}{A Deep Research Study — 1001 Examples, Pure Signal}}

\vspace{2.5cm}
{\sffamily\textcolor{accent}{Vision $\rightarrow$ Research $\rightarrow$ Mission}}\\[0.3em]
{\small\sffamily\textcolor{accent}{Full Pipeline Execution}}

\vspace{2.5cm}
{\small\sffamily\textcolor{subtlegray}{Date: March 26, 2026 \quad|\quad Status: Ready for GSD Orchestrator}}\\[0.3em]
{\small\sffamily\textcolor{subtlegray}{Activation Profile: Fleet (20 roles) \quad|\quad Est.: 18--24 context windows across 6--8 sessions}}\\[0.3em]
{\small\sffamily\textcolor{subtlegray}{Modules: 6 \quad|\quad Waves: 4 \quad|\quad Parallel Tracks: 3}}
\end{center}
\newpage

% ----------------------------------------------------------------
% TABLE OF CONTENTS
% ----------------------------------------------------------------
\tableofcontents
\newpage

% ================================================================
% STAGE 1 — VISION DOCUMENT
% ================================================================

\stageheader{STAGE 1 --- VISION DOCUMENT}{primary}

\section{GPU Graphics, Compute \& GSD-OS Observability Ecosystem --- Vision Guide}

\metafield{Date:}{March 26, 2026}
\metafield{Status:}{Vision Complete / Research Harvested / Ready for Execution}
\metafield{Depends On:}{GSD-OS Desktop Vision, GSD Chipset Architecture, Amiga New Workbench Vision}
\metafield{Context Windows:}{18--24 across 6--8 sessions}
\metafield{Activation Profile:}{Fleet (20 roles)}

\subsection{Vision}

The Amiga 500 shipped with a custom chipset that embarrassed workstations costing ten times its price. It did not win by having more transistors. It won because Agnus, Denise, and Paula each owned their slice of the machine and communicated through DMA channels that kept the CPU free to think. The result was a computer that could play four-channel audio, animate 32-color sprites, and receive input simultaneously at 7 MHz — a feat that would have required three dedicated processors in any conventional architecture of the era. This is the Amiga Principle made silicon: architectural leverage over brute force.

The GSD-OS desktop is an application of that same principle at the software layer. The CRT shader engine is the Denise chip — it owns the display and renders pixel-perfect phosphor bloom without asking permission. The skill-creator is Agnus — it watches memory usage patterns and pre-stages the next wave before the current one lands. The observability stack is Paula — it listens to everything, reports I/O health, and knows when a channel is saturated before the system does. The mission documented here exists to build the \emph{complete wiring harness} connecting these layers: from the sub-pixel OpenGL fragment shader to the Prometheus scrape interval, from the QPU assembler kernel on a Raspberry Pi 5 to the Grafana dashboard alert that fires when GPU temperature exceeds the ambient thermal envelope of a PNW server room.

The 1001 examples taxonomy is not hyperbole. The NeHe legacy produced 48 canonical OpenGL lessons that taught a generation. LearnOpenGL.com extended that to 60+ modern core-profile chapters. The Khronos Vulkan samples repository contains another 100+ production-quality demonstrations. Shadertoy alone hosts hundreds of thousands of GLSL fragment shaders, of which XScreenSaver 6.14 has begun pulling the open-licensed subset into its 240+-hack ecosystem. The GSD mission is to create a \emph{curated, annotated, cross-referenced} library of 1001 examples organized not as a flat list but as a unit circle of complexity — from the zero-degree ``Hello Triangle'' through the full rotation of ray marching, compute pipelines, and real-time telemetry rendering.

The Amiga expansion principles thread through the entire architecture as philosophy made hardware. Zorro II's autoconfig protocol — where expansion cards announce their own memory maps at boot — is conceptually identical to GSD skill-creator's autoload system, where skills register themselves into the active context on demand. The trapdoor slot that passed 68000 bus signals through to an accelerator card is the same idea as GSD's DACP three-part bundle: intent + data + executable, a structured handoff where the receiver knows exactly what to do without guessing. Every design choice in this mission should ask: \emph{does this card slot cleanly into the bus, or does it require a firmware patch to work?}

\subsection{Problem Statement}

\begin{enumerate}
  \item \textbf{The NeHe gap.} NeHe's 48-lesson curriculum was written for OpenGL 1.x immediate mode (1997--2012). LearnOpenGL.com covers modern core-profile but stops at OpenGL 4.3 practices, missing compute shaders, mesh shaders, ray tracing extensions, and Vulkan integration. No single resource unifies NeHe's pedagogical clarity with Vulkan's explicit GPU control model. Developers transitioning from tutorials to production encounter a cliff.

  \item \textbf{CUDA lock-in on non-NVIDIA silicon.} The Raspberry Pi 5 ships with a VideoCore VII GPU that supports OpenGL ES 3.1 and Vulkan 1.2, yet has no CUDA support and limited OpenCL coverage. Audio DSP workloads that could exploit parallel SIMD lanes on the QPU instead run on a single ARM Cortex-A76 core because no practical guide maps CUDA patterns to equivalent OpenCL/Vulkan compute or ARM NEON intrinsics. The SoC audio integration story is fragmented across five incompatible documentation ecosystems.

  \item \textbf{Amiga expansion knowledge is dying.} The Zorro bus specification, DMA arbitration protocol, autoconfig ROM mechanism, and chipset coprocessor design patterns that powered the 1985--1994 Amiga expansion ecosystem exist primarily in the Hardware Reference Manual, scattered forum posts, and the memories of engineers who are retiring. The homebrew community building FPGA accelerators (Vampire), new Zorro cards (AmigatronicZ2, Matthias Heinrichs ZorroLAN), and MiSTer FPGA cores is re-deriving this knowledge expensively. A structured reference would halve their iteration cycles.

  \item \textbf{XScreenSaver as unexploited GSD-OS render canvas.} XScreenSaver 6.14 contains 240+ OpenGL/GLSL screensaver hacks, a Shadertoy-compatible API wrapper, and a virtual root window system that makes the entire desktop a render target. GSD-OS's Tauri shell has no integration with this ecosystem. The CRT shader engine uses a WebGL pipeline that is structurally identical to XScreenSaver's GLSL wrapper but shares no code, documentation, or example library. This is a duplication of effort across two systems that want to do the same thing.

  \item \textbf{Observability dark spots in GPU workloads.} OpenTelemetry reached GA status for all signals in 2025. Prometheus 3.0 introduced resource attribute promotion for OTel metrics. NVIDIA DCGM Exporter surfaces GPU utilization, temperature, memory, and power as Prometheus metrics. Yet no unified GSD mission documents the wiring from a Raspberry Pi VideoCore temperature sensor through an OTel Collector configuration to a Grafana Alloy pipeline and into a GSD dashboard panel. The hardware exists; the plumbing specification does not.
\end{enumerate}

\subsection{Core Concept}

\textbf{The GPU as a specialized DMA coprocessor.} Just as Agnus owned chip RAM and the blitter owned the pixel pipeline, the GPU owns the framebuffer and the shader execution units. Every architecture decision in this mission follows from treating the GPU not as ``the CPU but parallel'' but as a \emph{dedicated co-processor with its own instruction set, memory bus, and DMA engine}. The CPU submits work orders; the GPU executes them asynchronously; the observability stack reads the health telemetry without interrupting either.

\subsubsection{Domain Architecture Map}

\begin{asciibox}
GSD-OS Desktop (Tauri Shell)
  |
  +-- Blueprint Editor (Canvas) -----------> WebGL / OpenGL ES 3.1
  |     |
  |     +-- CRT Shader Engine (Denise) ----> GLSL Fragment Shaders
  |     +-- Skill-Creator Feed (Agnus) ----> Live Metrics → Grafana Panel
  |
  +-- XScreenSaver Hacks (idle render) ---> GLX Virtual Root Window
  |     |
  |     +-- Shadertoy-compatible API -------> 240+ GLSL Hacks
  |     +-- Really Slick / OpenGL extras --> GLX Port
  |
  +-- Compute Layer
  |     |
  |     +-- Desktop GPU (NVIDIA/AMD) ------> CUDA / Vulkan Compute
  |     +-- Raspberry Pi 5 (VideoCore VII) -> OpenCL (VC4CL) / QPU ASM
  |     +-- Audio SoC DSP pipeline -------> ARM NEON / OpenCL kernels
  |
  +-- Observability Pipeline (Paula)
        |
        +-- NVIDIA DCGM Exporter ----------> Prometheus scrape endpoint
        +-- OTel Collector (Alloy) --------> OTLP → Prometheus → Grafana
        +-- Raspberry Pi hw metrics -------> node_exporter + custom
        +-- GSD skill-creator metrics -----> custom OTel SDK instrument
\end{asciibox}

\subsubsection{Module Architecture}

\begin{asciibox}
MODULE 01: OpenGL/Vulkan Foundations (1001 Examples Taxonomy)
  -> Hello Triangle → Core Profile → Compute Shaders → Vulkan → SPIR-V

MODULE 02: Compute Without NVIDIA (CUDA → OpenCL/Vulkan/NEON)
  -> CUDA concepts → OpenCL equivalents → Vulkan Compute → Pi QPU → Audio DSP

MODULE 03: Amiga Expansion Principles (Zorro Bus + Chipset Architecture)
  -> Zorro II/III autoconfig → DMA arbitration → FPGA/Vampire → GSD mapping

MODULE 04: XScreenSaver Engineering (GLX + GLSL + Shadertoy)
  -> virtual root window → GLSL portability → 240+ hack catalog → GSD-OS hooks

MODULE 05: GSD-OS Deep Integration
  -> Tauri shell → Blueprint Editor → WebGL CRT → skill-creator canvas

MODULE 06: Observability Stack (OTel + Prometheus + Grafana)
  -> DCGM exporter → OTel Collector → Alloy → Grafana dashboards → alerts
\end{asciibox}

\subsection{Research Modules}

\subsubsection{Module 01 — OpenGL/Vulkan Foundations: 1001 Examples Taxonomy}

The canonical learning ladder: NeHe's 48 lessons (immediate mode, historical context) → LearnOpenGL.com's 60+ modern core-profile chapters → open.gl (Alexander Overvoorde) → The Official Vulkan Tutorial (vulkan-tutorial.com) → vkguide.dev (dynamic rendering focus) → Khronos Official Vulkan Samples (100+) → Sascha Willems' Vulkan samples → Shadertoy GLSL library. This module catalogs all 1001 examples into a unit circle taxonomy: $\theta = 0$ (Hello Triangle) through $\theta = 2\pi$ (deferred rendering with ray tracing extensions). Each example is tagged with: API level (GL 2.1 / GL 4.6 / Vulkan 1.3), technique category, GSD-OS applicability, and Pi-5 compatibility.

\subsubsection{Module 02 — Compute Without NVIDIA: CUDA to OpenCL/Vulkan/NEON}

CUDA parallel programming patterns translated to their cross-vendor equivalents. The Raspberry Pi 5 VideoCore VII supports Vulkan 1.2 and OpenGL ES 3.1; the older VideoCore IV has VC4CL implementing OpenCL 1.2 embedded profile (approx.\ 24 GFLOPS). ARM NEON SIMD intrinsics provide 4× speedup on quad-core A76 without any GPU involvement. Audio DSP workloads (FFT, convolution reverb, pitch shifting) map naturally to OpenCL kernels or NEON vectorized C. This module documents the translation table from CUDA kernel patterns to OpenCL C / Vulkan Compute GLSL / ARM NEON intrinsics.

\subsubsection{Module 03 — Amiga Expansion Principles: Zorro Bus and Chipset Architecture}

The Zorro II bus (Amiga 2000) is a buffered 16-bit extension of the 68000 bus with DMA bus mastering. Zorro III (Amiga 3000/4000) extends to 32-bit with full address decode. The autoconfig mechanism allows cards to self-describe their memory requirements on boot — a 1985 plug-and-play architecture. The three custom chips (Agnus/Denise/Paula) each own dedicated DMA channels, executing concurrently with zero CPU overhead. Modern FPGA accelerators (Apollo/Vampire, MiSTer) re-implement these specs in programmable silicon. This module extracts the design principles that map directly onto GSD's chipset YAML (Agnus=orchestration, Denise=display, Paula=I/O) and DACP structured handoff architecture.

\subsubsection{Module 04 — XScreenSaver Engineering: GLX, GLSL, and Shadertoy}

XScreenSaver 6.14 (released January 2026) includes 18 new Shadertoy-derived hacks, a Shadertoy-compatible GLSL wrapper (``jwzgles'' — OpenGL 1.3 over OpenGL ES, ``jwxyz'' — X11 over Cocoa), and a virtual root window system where any X11 program can become a screensaver. The hack framework provides windowing, command-line parameter parsing, and preference saving, letting developers focus entirely on the render loop. The Really Slick Screensavers (RSS) GLX port provides another 10 high-quality OpenGL hacks. The GSD-OS CRT shader engine uses a structurally identical GLSL pipeline. This module documents the XScreenSaver hack API, GLSL portability matrix (GL 2.1 / GL 4.6 / ES 3.0 / WebGL 2.0), and integration path into GSD-OS as an idle-state render mode.

\subsubsection{Module 05 — GSD-OS Deep Integration}

The Tauri-based GSD-OS desktop shell connects the Blueprint Editor (canvas), the skill-creator feed, and the CRT shader engine. This module documents: the WebGL pipeline inside the Tauri WebView, the Bridge between Tauri Rust core and the JS render loop, the event bus that carries skill-creator metrics to the canvas, and the hook points for injecting XScreenSaver-compatible GLSL programs as idle-state wallpapers. The GSD chipset YAML (Agnus/Denise/Paula/68000/Gary) is formalized here as a first-class configuration artifact with version-controlled defaults.

\subsubsection{Module 06 — Observability Stack: OpenTelemetry, Prometheus, and Grafana}

Full-stack GPU observability from hardware sensor to dashboard alert. NVIDIA DCGM Exporter surfaces GPU metrics (utilization, temperature, power, memory) in Prometheus format. VC4CL and VideoCore VII on Pi expose temperature and utilization via sysfs. The OTel Collector (Grafana Alloy, announced GrafanaCON 2024, Grafana Agent EOL November 2025) pipelines all metrics with OTLP, leveraging Prometheus 3.0's resource attribute promotion (GA, Grafana Cloud default). GSD skill-creator is instrumented with the OTel SDK to emit custom spans. Final output: 5 pre-built Grafana dashboards covering GPU health, Raspberry Pi SoC, audio DSP latency, GSD skill metrics, and the full observability pipeline health.

\subsection{Chipset Configuration}

\begin{asciibox}
# GSD Chipset YAML — GPU/GSD-OS/Observability Mission
chipset:
  agnus:                          # Orchestration / DMA / Context
    model: claude-opus-4-6
    role: FLIGHT + INTEG + ARCH
    responsibilities:
      - Cross-module synthesis and dependency graph management
      - Amiga expansion principle extraction and GSD mapping
      - Wave gate approval; context budget management
    allocation: 30%

  denise:                         # Display / Output Rendering
    model: claude-sonnet-4-6
    role: EXEC-1 + EXEC-2 + CRAFT-GL + CRAFT-OBS
    responsibilities:
      - OpenGL/Vulkan examples taxonomy compilation (Module 01)
      - XScreenSaver hack catalog and GLSL portability matrix (Module 04)
      - Observability stack documentation (Module 06)
      - GSD-OS WebGL pipeline specification (Module 05)
    allocation: 60%

  paula:                          # Audio / I/O / Anticipatory
    model: claude-haiku-4-5
    role: SCOUT + LOG + BUDGET
    responsibilities:
      - Source index and shared schema generation (Wave 0)
      - Token budget tracking and session health telemetry
      - Commit messages and milestone summaries
    allocation: 10%

\end{asciibox}

\subsection{Scope Boundaries}

\subsubsection{In Scope (v1.0)}

\begin{itemize}[noitemsep]
  \item OpenGL 2.1 through 4.6 core profile examples (annotated NeHe modernizations + LearnOpenGL curriculum)
  \item Vulkan 1.2/1.3 introduction through deferred rendering pipeline
  \item GLSL / SPIR-V shading language reference with Shadertoy compatibility notes
  \item CUDA-to-OpenCL/Vulkan-Compute pattern translation table
  \item Raspberry Pi VideoCore VII (Pi 5): Vulkan 1.2, OpenGL ES 3.1, VC4CL (Pi 3/4)
  \item ARM NEON SIMD audio DSP patterns (FFT, convolution, pitch)
  \item Amiga Zorro II/III bus specification and autoconfig mechanism
  \item Amiga custom chipset DMA architecture (Agnus/Denise/Paula coprocessor model)
  \item Modern Amiga FPGA accelerators: Apollo/Vampire, MiSTer, AmigatronicZ2, ZorroLAN-IDE
  \item XScreenSaver 6.14 hack framework API (GLX virtual root window, GLSL wrapper)
  \item Shadertoy-compatible module integration for GSD-OS idle render
  \item Tauri shell → WebGL bridge hook points
  \item GSD chipset YAML formalization and version-controlled defaults
  \item NVIDIA DCGM Exporter → Prometheus scrape configuration
  \item Grafana Alloy OTel Collector pipeline (OTLP + Prometheus resource attribute promotion)
  \item 5 pre-built Grafana dashboard specifications
  \item GSD skill-creator OTel SDK instrumentation hooks
\end{itemize}

\subsubsection{Out of Scope}

\begin{itemize}[noitemsep]
  \item DirectX 12 / Metal / WebGPU (separate mission)
  \item NVIDIA ray tracing hardware (RTX) deep dive — referenced, not primary
  \item AmigaOS 3.x/4.x application development
  \item Machine learning / AI inference on GPU (separate GSD mission)
  \item Production Kubernetes deployment of the observability stack
  \item Windows-specific screensaver (.scr) API
\end{itemize}

\subsection{Success Criteria}

\begin{enumerate}
  \item A 1001-example taxonomy is compiled, tagged, and cross-referenced with API level, technique category, Pi-5 compatibility, and GSD-OS applicability. Every example has a source URL or repository reference.
  \item The NeHe gap is bridged: all 48 NeHe lessons have a modern GL 4.x or Vulkan equivalent documented with GLSL/SPIR-V shader source.
  \item The CUDA-to-OpenCL translation table covers at minimum: parallel reduction, FFT, matrix multiply, convolution, histogram, prefix scan. Each entry specifies API, kernel signature pattern, and performance characteristics on Pi 5.
  \item Raspberry Pi 5 audio DSP pipeline is documented end-to-end: ALSA/JACK input → ARM NEON or OpenCL kernel → Vulkan compute → audio output with measured latency in milliseconds.
  \item Amiga Zorro II/III autoconfig protocol is fully documented with a GSD architectural mapping table (Zorro concept → GSD DACP equivalent).
  \item Amiga chipset DMA coprocessor model (Agnus/Denise/Paula) is mapped 1:1 onto GSD chipset YAML roles with rationale for each assignment.
  \item XScreenSaver hack framework API is documented with a working Hello-Triangle equivalent hack skeleton and GLSL portability matrix across GL 2.1, GL 4.6, ES 3.0, and WebGL 2.0.
  \item A curated 50-hack subset of XScreenSaver/Shadertoy is cataloged for GSD-OS idle render, each annotated for aesthetic category and GPU load.
  \item GSD-OS Tauri integration points are specified: WebGL context creation, GLSL program injection, skill-creator metric event bus, and idle-state render mode activation.
  \item NVIDIA DCGM Exporter + OTel Collector (Grafana Alloy) pipeline is documented with complete YAML configuration from scrape endpoint to Prometheus remote write.
  \item Raspberry Pi SoC temperature and VideoCore utilization metrics are integrated into the same OTel pipeline.
  \item Five Grafana dashboard specifications are produced: GPU Health, Pi SoC, Audio DSP Latency, GSD Skill Metrics, and Pipeline Health.
\end{enumerate}

\subsection{Relationship to Other GSD Documents}

\begin{center}
\rowcolors{2}{rowgray}{white}
\begin{tabularx}{\textwidth}{L{4.5cm}XC{2.5cm}}
\toprule
\rowcolor{primary}
\tableheader{Document} & \tableheader{Relationship} & \tableheader{Dependency} \\
\midrule
GSD-OS Desktop Vision & Parent vision; defines Tauri shell, CRT shader engine & Upstream \\
GSD Chipset Architecture Vision & Agnus/Denise/Paula model formalized here & Upstream \\
GSD Amiga New Workbench Vision & Workbench UI aesthetic depends on XScreenSaver idle renders & Downstream \\
GSD Staging Layer Vision & DACP bundles carry GLSL programs between agents & Peer \\
Deep Audio Analyzer Mission & Audio DSP pipeline (Module 02) directly informs DAA signal chain & Peer \\
Skill Creator Wasteland Integration & OTel instrumentation of skill-creator (Module 06) & Peer \\
Fox Infrastructure Group (FoxCompute) & Ocean compute nodes will run the Vulkan Compute workloads & Downstream \\
\bottomrule
\end{tabularx}
\end{center}

\subsection{The Through-Line}

The Amiga shipped with a chassis that looked like a toy and a chipset that looked like a miracle. The miracle was not the hardware. The miracle was the decision to give each coprocessor its own lane of the bus and trust it to drive its own DMA without supervision. The CPU was freed to do what only the CPU could do: think. Everything else ran in the spaces between clock cycles.

GSD-OS is a software architecture built on the same decision. The shader engine owns the display lane. The skill-creator owns the context lane. The observability stack owns the telemetry lane. Each runs its own DMA equivalent — its own asynchronous event loop — without blocking the others. The 1001 examples in this mission are not a curriculum. They are the unit circle of what those lanes can carry: from a single shaded triangle at $\theta = 0$ to a full deferred renderer with live Grafana telemetry at $\theta = 2\pi$. The spaces between each example are where the real learning happens.

\newpage

% ================================================================
% STAGE 2 — RESEARCH REFERENCE
% ================================================================

\stageheader{STAGE 2 --- RESEARCH REFERENCE}{secondary}

\section{GPU Graphics, Compute \& GSD-OS Observability --- Research Reference}

\subsection{How to Use This Document}

This reference is organized by module. Each section contains sourced findings, quantitative data, and direct references to authoritative repositories. Agents executing Wave 1 should treat each module section as their primary source document.

\textbf{Key source organizations:}
\begin{itemize}[noitemsep]
  \item Khronos Group (khronos.org, vulkan.org, opengl.org) --- OpenGL and Vulkan specifications
  \item The Khronos Group Vulkan Samples Repository (github.com/KhronosGroup)
  \item LearnOpenGL.com (Joey de Vries) --- Modern OpenGL core-profile curriculum
  \item vulkan-tutorial.com --- Canonical Vulkan introduction
  \item vkguide.dev --- Dynamic rendering-focused Vulkan guide
  \item XScreenSaver (jwz.org) --- Jamie Zawinski's screensaver ecosystem
  \item Raspberry Pi Foundation (raspberrypi.com) --- VideoCore GPU documentation
  \item NVIDIA DCGM (docs.nvidia.com/datacenter/dcgm) --- GPU metric exporters
  \item Grafana Labs (grafana.com) --- Alloy, Prometheus, Loki, Tempo
  \item AMD GPUOpen (gpuopen.com) --- Vulkan best practices and profiling
\end{itemize}

\subsection{Module 01 Reference --- OpenGL/Vulkan Foundations}

\subsubsection{The NeHe Legacy and the Modern Gap}

NeHe Productions (nehe.gamedev.net, 1997--2012) published 48 lessons covering OpenGL 1.x immediate mode. The site remains archived and widely referenced as pedagogical history. The core gap: immediate-mode functions (glBegin, glEnd, glLight) are deprecated in OpenGL 3.2 core profile and unavailable in OpenGL ES 3.x or WebGL.

LearnOpenGL.com (learnopengl.com) fills this gap with 60+ modern core-profile chapters, available as both a free website and a commercial print book. The curriculum covers: Getting Started → Lighting → Model Loading → Advanced OpenGL → Advanced Lighting → PBR → In Practice → Guest Articles. The site has been continuously revised over 7 years.

The Vulkan Tutorial (vulkan-tutorial.com) by Alexander Overvoorde is the canonical Vulkan introduction, requiring C++ with RAII familiarity and C++17 compiler support (GCC 7+, Clang 5+, MSVC 2017+). It explicitly does not assume OpenGL knowledge but does require 3D graphics fundamentals.

\textbf{1001 Examples Source Breakdown:}

\begin{center}
\rowcolors{2}{rowgray}{white}
\begin{tabularx}{\textwidth}{L{4.5cm}C{1.8cm}XC{2cm}}
\toprule
\rowcolor{primary}
\tableheader{Source} & \tableheader{Count} & \tableheader{API Level} & \tableheader{Pi-5 Compat} \\
\midrule
NeHe Productions (modernized) & 48 & GL 2.1--4.x & Partial \\
LearnOpenGL.com chapters & 65 & GL 3.3--4.6 & ES 3.1 subset \\
open.gl (Overvoorde) & 30 & GL 3.x & Yes (ES 3.1) \\
vulkan-tutorial.com & 25 & Vulkan 1.2+ & Yes (VK 1.2) \\
vkguide.dev & 30 & Vulkan 1.3 & Yes (VK 1.2) \\
Khronos Official Vulkan Samples & 120 & Vulkan 1.x & Partial \\
Sascha Willems Vulkan Samples & 140 & Vulkan 1.x & Partial \\
XScreenSaver GL Hacks & 80 & GL 1.3--4.x & GL ES 3.0 \\
Shadertoy (open-licensed) & 200 & GLSL / WebGL & WebGL 2.0 \\
Really Slick Screensavers (GLX) & 10 & GL 1.x--2.x & Partial \\
Custom GSD-OS examples & 253 & Mixed & Yes \\
\textbf{Total} & \textbf{1001} & --- & --- \\
\bottomrule
\end{tabularx}
\end{center}

\subsubsection{GLSL Version Portability Matrix}

XScreenSaver 6.14 author Jamie Zawinski documented the portability challenge directly: targeting four incompatible GLSL versions — Linux GL 3.1--4.6 compatibility profile, macOS GL 4.1 core profile, iOS OpenGL ES 3.0, and Android OpenGL ES 3.0--3.2. He implemented compatibility shims to bridge GL ES 3.0 code to GL 2.1 systems.

\begin{center}
\rowcolors{2}{rowgray}{white}
\begin{tabularx}{\textwidth}{L{3cm}C{2cm}C{2cm}C{2.5cm}C{2.5cm}}
\toprule
\rowcolor{primary}
\tableheader{Feature} & \tableheader{GL 2.1} & \tableheader{GL 4.6} & \tableheader{ES 3.0/WebGL2} & \tableheader{Vulkan SPIR-V} \\
\midrule
Compute Shaders & No & Yes & No & Yes \\
Geometry Shaders & No & Yes & No & Yes \\
Tessellation & No & Yes & No & Yes \\
Integer Types & Limited & Yes & Yes & Yes \\
Mesh Shaders & No & Ext & No & Ext \\
Ray Tracing & No & Ext & No & Ext \\
Bindless Textures & No & Ext & No & Yes \\
\bottomrule
\end{tabularx}
\end{center}

\subsection{Module 02 Reference --- Compute Without NVIDIA}

\subsubsection{Raspberry Pi GPU Landscape}

The Raspberry Pi 3 Model B+ uses a Broadcom BCM2837B0 with VideoCore IV GPU. VC4CL (VC for Compute Language), implemented by Daniel Steadelmann at Nuremberg Tech, provides OpenCL 1.2 embedded profile support. The VideoCore IV's maximum compute performance is approximately 24 GFLOPS — insufficient for neural network training but adequate for real-time DSP workloads.

The Raspberry Pi 5 ships with Broadcom BCM2712 and VideoCore VII, supporting Vulkan 1.2 natively. Neither VideoCore generation supports CUDA (NVIDIA proprietary) or the full OpenCL specification.

The critical alternative path: ARM NEON SIMD intrinsics via \texttt{\#pragma omp parallel for} provide 4× speedup on quad-core A76 processors without GPU involvement. For audio DSP workloads, NEON vectorized FFT (e.g., KissFFT, FFTW) often outperforms VideoCore IV OpenCL due to DMA transfer overhead.

\subsubsection{CUDA to Cross-Vendor Translation}

\begin{center}
\rowcolors{2}{rowgray}{white}
\begin{tabularx}{\textwidth}{L{3cm}L{3cm}L{3cm}X}
\toprule
\rowcolor{primary}
\tableheader{CUDA Concept} & \tableheader{OpenCL Equiv.} & \tableheader{Vulkan Compute} & \tableheader{Pi-5 / NEON} \\
\midrule
cudaMalloc & clCreateBuffer & VkBuffer + VkDeviceMemory & malloc + posix\_memalign \\
cudaMemcpy & clEnqueueCopy & vkCmdCopyBuffer & memcpy / DMA \\
\_\_global\_\_ kernel & \_\_kernel function & layout(local\_size) comp shader & NEON intrinsic loop \\
blockDim / gridDim & get\_global\_id() & gl\_GlobalInvocationID & omp parallel for \\
cudaDeviceSynchronize & clFinish & vkQueueWaitIdle & pthread\_join \\
shared memory & \_\_local memory & shared (workgroup storage) & stack-allocated array \\
warp (32 threads) & wavefront (64/32) & subgroup & NEON vector lane \\
\bottomrule
\end{tabularx}
\end{center}

SYCL (C++ abstraction over OpenCL) and Intel oneAPI/DPC++ represent the current state-of-the-art in CUDA portability. ROCm/HIP translates CUDA code to run on AMD GPUs. For embedded SoC work, Vulkan Compute provides the most consistent path across VideoCore VII, Mali, and Adreno GPUs.

\subsubsection{Audio DSP on SoC}

Key audio processing workloads and their vectorization strategies on Pi 5:

\begin{itemize}[noitemsep]
  \item \textbf{FFT (Fast Fourier Transform):} FFTW with NEON backend; alternatively Vulkan Compute FFT via VkFFT library
  \item \textbf{Convolution Reverb:} Frequency-domain multiplication (FFT $\times$ IR); NEON float32x4 multiply-accumulate
  \item \textbf{Pitch Shifting:} Phase vocoder on FFT bins; parallelizes cleanly across NEON lanes
  \item \textbf{Parametric EQ / IIR filters:} Direct Form II transposed; 4-sample NEON pipeline
  \item \textbf{Granular Synthesis:} Embarrassingly parallel grain rendering; ideal Vulkan Compute candidate
\end{itemize}

Measured latency target for real-time audio on Pi 5: $<$ 10ms round-trip with JACK audio server, 128-frame buffer at 44100 Hz.

\subsection{Module 03 Reference --- Amiga Expansion Principles}

\subsubsection{Zorro Bus Architecture}

Zorro II (Amiga 2000) is a buffered 16-bit extension of the Motorola 68000 bus with DMA bus mastering support. Zorro III (Amiga 3000/4000) extends to 32-bit addressing and full address decode. The Amiga 4000 provides four Zorro III slots. Fast RAM can be extended via Zorro III cards to 64, 128, or 256 MiB.

The autoconfig mechanism is the architectural gem: each expansion card carries a ROM that announces its memory requirements, interrupt vectors, and device ID to the system on first boot. The OS allocates address space dynamically. No manual jumper configuration required. This is functionally identical to modern PCI-E enumeration — implemented in 1985.

\subsubsection{Chipset Coprocessor Architecture}

The three Amiga custom chips represent the canonical specialized-coprocessor model:

\begin{center}
\rowcolors{2}{rowgray}{white}
\begin{tabularx}{\textwidth}{L{2cm}L{3.5cm}XL{3cm}}
\toprule
\rowcolor{primary}
\tableheader{Chip} & \tableheader{Amiga Role} & \tableheader{Capabilities} & \tableheader{GSD YAML Role} \\
\midrule
Agnus & DMA controller / Blitter & Owns chip RAM, runs blitter/copper independently & FLIGHT / Orchestration \\
Denise & Display / Color & Sprites, playfields, color register, genlock & Display / CRT shader \\
Paula & Audio / I/O & 4-channel PCM, disk DMA, serial/parallel I/O & Audio / I/O / Scout \\
Gary & Glue logic & Address decode, bus arbitration, peripheral enables & 68000 / Glue \\
\bottomrule
\end{tabularx}
\end{center}

\subsubsection{Modern FPGA Accelerators}

The Apollo/Vampire project provides FPGA-based Amiga accelerators with CPU performance beyond the 68060, SAGA graphics at Full-HD resolution, HDMI output, 128 MB RAM, and MicroSD storage. Individual Computers (iComp) produces the iSprint accelerator (MC68EC000 at 14--42 MHz, 8 MB RAM, 8 MB FlashROM). AmigatronicZ2 (Italy) manufactures a Zorro II side-expansion card enabling Amiga 500 owners to use full Zorro II cards. A custom Zorro card by Matthias Heinrichs combining IDE, clock port, and Ethernet has been successfully assembled and benchmarked at 2.8--3 MB/s IDE transfer rates.

\subsection{Module 04 Reference --- XScreenSaver Engineering}

\subsubsection{Framework Architecture}

XScreenSaver (Jamie Zawinski, 1992--present) is a free and open-source collection of 240+ screensavers for Unix, macOS, iOS, and Android. Version 6.14 (January 2026) added 18 Shadertoy-derived hacks via a new Shadertoy-compatible API wrapper. The framework includes:

\begin{itemize}[noitemsep]
  \item \textbf{jwxyz:} Complete X11 API implementation over Cocoa (macOS/iOS)
  \item \textbf{jwzgles:} OpenGL 1.3 API over OpenGL ES 1.0 (iOS/Android)
  \item \textbf{Virtual root window:} Hacks draw to the virtual root window rather than a normal window, enabling them to cover the entire desktop
  \item \textbf{GLX integration:} OpenGL hacks use GLX (not GLFW) to hook into the virtual root window; GLFW is incompatible with the screensaver framework
\end{itemize}

Writing a new hack requires: including \texttt{vroot.h}, using GLX for context creation, implementing the screensaver protocol (handle expose events, periodic redraws, respond to SIGTERM). The existing hack source code in the XScreenSaver distribution provides complete examples.

\subsubsection{Shadertoy Licensing Note}

Shadertoy's default license is CC-BY-NC-SA (non-commercial, share-alike). This is incompatible with XScreenSaver's open-source license. Only shaders licensed MIT, BSD, CC-BY, CC-BY-SA, or CC0 can be redistributed with XScreenSaver. This licensing constraint shapes the curated GSD-OS 50-hack subset.

\subsubsection{50-Hack Curated Subset for GSD-OS}

Categories for the GSD-OS idle render catalog:

\begin{center}
\rowcolors{2}{rowgray}{white}
\begin{tabularx}{\textwidth}{L{4cm}C{1.5cm}XC{2cm}}
\toprule
\rowcolor{primary}
\tableheader{Category} & \tableheader{Count} & \tableheader{Representative Hacks} & \tableheader{GPU Load} \\
\midrule
Fractal / Mathematical & 8 & glmatrix, hypercube, polytopes, glknots & Low--Med \\
Particle Systems & 7 & xrayswarm, nerverot, fiberlamp & Low--Med \\
3D Geometry & 9 & flipflop, blocktube, topblock, cubestorm & Medium \\
Amiga-Nostalgic & 5 & boing, phosphor, starwars, atunnel & Low \\
Shadertoy GLSL & 12 & alien beacon, battered planet (new in 6.14) & Med--High \\
Really Slick Screensavers & 5 & Skyrocket, Helios, Euphoria, Flux & Medium \\
Generative Art & 4 & glslideshow, cloudlife, gflux, xflame & Low \\
\bottomrule
\end{tabularx}
\end{center}

\subsection{Module 05 Reference --- GSD-OS Deep Integration}

Key Tauri integration points for the WebGL/GLSL pipeline:

\begin{itemize}[noitemsep]
  \item \textbf{WebView GLSL context:} Tauri's WebView uses the system WebGL2 implementation. On Linux: Mesa with GL 4.x compatibility; on Pi 5: WebGL 2.0 over VideoCore VII
  \item \textbf{GLSL injection path:} Blueprint Editor canvas exposes a \texttt{registerShaderProgram(id, vertSrc, fragSrc)} API callable from Rust core via Tauri commands
  \item \textbf{Skill-creator metric bus:} Tauri event system (\texttt{emit}/\texttt{listen}) carries skill-creator YAML diffs to the canvas as structured JSON; the canvas renders them as animated overlay panels
  \item \textbf{Idle-state render:} When GSD-OS detects $>$ 5 minutes of inactivity, it fires a \texttt{screensaver:activate} event; the canvas switches to the selected GLSL hack at full-screen WebGL
  \item \textbf{CRT shader engine:} The existing phosphor/scanline/bloom fragment shader is the Denise chip of GSD-OS; it runs as a post-processing pass over all canvas output
\end{itemize}

\subsection{Module 06 Reference --- Observability Stack}

\subsubsection{NVIDIA DCGM Exporter Pipeline}

NVIDIA's Data Center GPU Manager (DCGM) suite exposes GPU metrics in Prometheus format via dcgm-exporter. Key metrics: \texttt{DCGM\_FI\_DEV\_GPU\_UTIL} (utilization \%), \texttt{DCGM\_FI\_DEV\_GPU\_TEMP} (temperature °C), \texttt{DCGM\_FI\_DEV\_POWER\_USAGE} (watts), \texttt{DCGM\_FI\_DEV\_MEM\_COPY\_UTIL} (memory bandwidth \%).

Grafana Cloud's AI Observability integration (GA, February 2025) instruments over 50 gen AI tools via OpenLIT (OpenTelemetry-native SDK) and provides pre-built GPU dashboards tracking utilization, temperature, and power.

\subsubsection{Grafana Alloy and Prometheus 3.0}

Grafana Alloy (announced GrafanaCON 2024) is the successor to Grafana Agent (EOL November 1, 2025). It is 100\% OTLP-compatible and provides native pipelines for both OpenTelemetry and Prometheus formats. Prometheus 3.0 introduced resource attribute promotion: OTel resource attributes (service.name, service.instance.id, k8s.pod.name) are promoted as Prometheus metric labels, eliminating the need for complex \texttt{target\_info} joins in PromQL queries.

According to Grafana Labs' 2024 Observability Survey, 89\% of respondents invest in Prometheus and 85\% in OpenTelemetry, with nearly 40\% using both in production.

\subsubsection{OTel Collector Configuration Skeleton}

\begin{asciibox}
# Grafana Alloy / OTel Collector — GSD-OS GPU Pipeline
receivers:
  prometheus:
    config:
      scrape_configs:
        - job_name: 'dcgm-gpu-metrics'
          scrape_interval: 15s
          static_configs:
            - targets: ['localhost:9400']
          metric_relabel_configs:
            - source_labels: [__name__]
              regex: 'DCGM_.*'
              action: keep
        - job_name: 'pi5-soc-metrics'
          scrape_interval: 30s
          static_configs:
            - targets: ['pi5.local:9100']   # node_exporter
        - job_name: 'gsd-skill-creator'
          scrape_interval: 10s
          static_configs:
            - targets: ['localhost:8888']   # OTel SDK endpoint
  otlp:
    protocols:
      grpc:
        endpoint: 0.0.0.0:4317
      http:
        endpoint: 0.0.0.0:4318

processors:
  batch:
    timeout: 10s
    send_batch_size: 1024
  resource:
    attributes:
      - action: insert
        key: gsd.environment
        value: gsd-os-desktop

exporters:
  prometheusremotewrite:
    endpoint: "http://prometheus:9090/api/v1/write"
  otlphttp/grafana-cloud:
    endpoint: "https://otlp-gateway.grafana.net/otlp"
    headers:
      Authorization: "Basic ${GRAFANA_CLOUD_TOKEN}"

service:
  pipelines:
    metrics:
      receivers: [prometheus, otlp]
      processors: [batch, resource]
      exporters: [prometheusremotewrite, otlphttp/grafana-cloud]
\end{asciibox}

\subsection{Safety and Sensitivity Considerations}

\begin{center}
\rowcolors{2}{rowgray}{white}
\begin{tabularx}{\textwidth}{L{4cm}L{2.5cm}X}
\toprule
\rowcolor{primary}
\tableheader{Area} & \tableheader{Type} & \tableheader{Handling} \\
\midrule
Shader code examples & Technical accuracy & All GLSL/SPIR-V examples must compile without errors on target platform \\
Open-source licensing & Legal & Shadertoy integration requires CC0/MIT/BSD/CC-BY only; CC-BY-NC-SA explicitly excluded \\
Security (OTel tokens) & Security & API keys and cloud tokens documented as environment variables, never hardcoded \\
Raspberry Pi safety & Hardware & QPU assembly notes must include DMA MMU warning (root requirement, memory access risk) \\
Performance claims & Numerical attribution & All GFLOPS and latency figures attributed to specific benchmarks or datasheets \\
\bottomrule
\end{tabularx}
\end{center}

\subsection{Source Bibliography}

\textbf{Khronos Group / Official Specifications}
\begin{itemize}[noitemsep]
  \item Vulkan 1.3 Specification — docs.vulkan.org
  \item OpenGL 4.6 Core Profile Specification — opengl.org
  \item Khronos Vulkan Samples Repository — github.com/KhronosGroup/Vulkan-Samples
  \item Khronos Beginners Guide to Vulkan — khronos.org/blog (2017)
\end{itemize}

\textbf{Learning Resources (Peer-quality, open curriculum)}
\begin{itemize}[noitemsep]
  \item LearnOpenGL.com — Joey de Vries (7-year continuously revised curriculum)
  \item vulkan-tutorial.com — Alexander Overvoorde (canonical Vulkan introduction)
  \item vkguide.dev — Victor Blanco (dynamic rendering focus, active 2024)
  \item open.gl — Alexander Overvoorde (modern OpenGL, interactive examples)
\end{itemize}

\textbf{Official Platform Documentation}
\begin{itemize}[noitemsep]
  \item Raspberry Pi Documentation — raspberrypi.com/documentation
  \item Broadcom BCM2712 VideoCore VII Datasheet
  \item VC4CL: OpenCL for VideoCore IV — github.com/doe300/VC4CL (master thesis, Nuremberg Tech)
  \item NVIDIA DCGM Documentation — docs.nvidia.com/datacenter/dcgm
  \item Amiga Hardware Reference Manual — amigadev.elowar.com
\end{itemize}

\textbf{Observability Stack}
\begin{itemize}[noitemsep]
  \item Grafana Alloy — grafana.com/docs/alloy (GrafanaCON 2024 announcement)
  \item Prometheus 3.0 Resource Attribute Promotion — grafana.com/blog (2025)
  \item Grafana AI Observability with GPU Monitoring — grafana.com/whats-new (Feb 2025)
  \item OpenTelemetry Specification — opentelemetry.io/docs (GA all signals, 2025)
\end{itemize}

\textbf{XScreenSaver / Screensaver Engineering}
\begin{itemize}[noitemsep]
  \item XScreenSaver 6.14 Release Notes — jwz.org/blog (January 2026)
  \item XScreenSaver Source Repository — github.com/porridge/xscreensaver
  \item Writing X11 Screensavers with OpenGL — dbeef.lol (2018, GLX virtual root window)
  \item Really Slick Screensavers GLX Port — sourceforge.net/projects/rssavers
\end{itemize}

\textbf{Amiga Hardware}
\begin{itemize}[noitemsep]
  \item Amiga Hardware Database — amiga.resource.cx
  \item Amiga Zorro II — en.wikipedia.org/wiki/Amiga\_Zorro\_II
  \item Building an Amiga Zorro Network Card (GadgetUK, 2025) — theoasisbbs.com
  \item AMYA500Z2 Zorro II Side Expansion — amigatronic.com / Tindie (2024)
  \item Apollo/Vampire FPGA Accelerator — apollo-accelerators.com
\end{itemize}

\newpage

% ================================================================
% STAGE 3 — MISSION PACKAGE
% ================================================================

\stageheader{STAGE 3 --- MISSION PACKAGE}{tertiary}

\section{Milestone Specification}

\subsection{Mission Objective}

Produce a publication-quality reference suite covering OpenGL/Vulkan graphics programming (1001 annotated examples), cross-vendor GPU compute translation, Amiga expansion architecture principles, XScreenSaver/GLSL engineering, GSD-OS deep integration specifications, and a complete OpenTelemetry/Prometheus/Grafana observability stack — all wired together under the Amiga Principle of specialized coprocessors executing in parallel without blocking the CPU.

\subsection{Architecture Overview}

\begin{asciibox}
WAVE 0: Foundation (Sequential — Paula/Haiku)
  Source index, shared schemas, example taxonomy skeleton, YAML templates
  ↓
WAVE 1: Parallel Research (3 Tracks — Denise/Sonnet + Agnus/Opus)
  Track A: Modules 01+02 (OpenGL/Vulkan + Compute/Pi)     [Sonnet]
  Track B: Modules 03+04 (Amiga + XScreenSaver)           [Sonnet + Opus]
  Track C: Modules 05+06 (GSD-OS + Observability)         [Sonnet]
  ── CAPCOM GATE: scope confirm, licensing review ──
  ↓
WAVE 2: Synthesis (Sequential — Agnus/Opus)
  Cross-module integration; Amiga→GSD mapping tables
  Unit circle taxonomy finalization; GLSL portability matrix
  ── CAPCOM GATE: technical accuracy review ──
  ↓
WAVE 3: Publication + Verification (Sequential — Denise/Sonnet)
  LaTeX compilation, 5 Grafana dashboard specs, index.html
  OTel config YAML validation; bibliography audit
\end{asciibox}

\subsection{Deliverables}

\begin{center}
\rowcolors{2}{rowgray}{white}
\begin{tabularx}{\textwidth}{C{0.5cm}L{4.5cm}XC{2.5cm}}
\toprule
\rowcolor{primary}
\tableheader{\#} & \tableheader{Deliverable} & \tableheader{Acceptance Criteria} & \tableheader{Format} \\
\midrule
D1 & 1001 Examples Taxonomy & All 1001 entries tagged: API, technique, Pi-5 compat, GSD applicability & Markdown table + JSON \\
D2 & CUDA$\to$OpenCL Translation Table & 7+ CUDA patterns with OpenCL/Vulkan/NEON equivalents & Markdown table \\
D3 & Amiga$\to$GSD Mapping Document & Zorro autoconfig $\to$ DACP; chipset $\to$ chipset YAML & Markdown \\
D4 & XScreenSaver Hack Catalog & 50-hack curated subset with GLSL portability notes & Markdown table \\
D5 & GSD-OS Integration Spec & Tauri hook points, GLSL injection API, metric event bus & Markdown \\
D6 & OTel Collector YAML Config & Complete Alloy config, validated syntax & YAML \\
D7 & 5 Grafana Dashboard Specs & JSON import files or panel specifications & JSON \\
D8 & Hello-Triangle XScreenSaver Skeleton & Compiles and runs as an XScreenSaver GLX hack & C source \\
D9 & Research Reference PDF & Full three-stage document & PDF/LaTeX \\
\bottomrule
\end{tabularx}
\end{center}

\subsection{Component Breakdown}

\begin{center}
\rowcolors{2}{rowgray}{white}
\begin{tabularx}{\textwidth}{L{3.5cm}L{3cm}C{2cm}C{2cm}}
\toprule
\rowcolor{primary}
\tableheader{Component} & \tableheader{Dependencies} & \tableheader{Model} & \tableheader{Est. Tokens} \\
\midrule
Wave 0: Source Index & None & Haiku & 8K \\
Wave 0: Shared Schemas & None & Haiku & 5K \\
Wave 0: Example Taxonomy Shell & Source Index & Haiku & 10K \\
Module 01: GL/VK Foundations & Taxonomy Shell & Sonnet & 40K \\
Module 02: Compute/Pi & Source Index & Sonnet & 35K \\
Module 03: Amiga Architecture & Source Index & Opus & 30K \\
Module 04: XScreenSaver & Source Index & Sonnet & 25K \\
Module 05: GSD-OS Integration & Modules 01,04 & Sonnet & 30K \\
Module 06: Observability Stack & Source Index & Sonnet & 35K \\
Wave 2: Amiga$\to$GSD Synthesis & Modules 01--06 & Opus & 25K \\
Wave 2: Unit Circle Taxonomy & All modules & Opus & 20K \\
Wave 3: OTel YAML + Dashboards & Module 06 & Sonnet & 20K \\
Wave 3: Bibliography Audit & All modules & Sonnet & 10K \\
Wave 3: Final Compilation & All & Sonnet & 15K \\
\textbf{Total (est.)} & & & \textbf{308K} \\
\bottomrule
\end{tabularx}
\end{center}

\subsection{Model Assignment Rationale}

\begin{center}
\rowcolors{2}{rowgray}{white}
\begin{tabularx}{\textwidth}{L{2cm}C{2cm}C{1.5cm}X}
\toprule
\rowcolor{primary}
\tableheader{Role} & \tableheader{Model} & \tableheader{Budget \%} & \tableheader{Rationale} \\
\midrule
Opus (Agnus) & claude-opus-4-6 & 30\% & Synthesis, Amiga architecture judgment, cross-module integration, DACP mapping decisions \\
Sonnet (Denise) & claude-sonnet-4-6 & 60\% & All six module documents, YAML configs, dashboard specs, compilation \\
Haiku (Paula) & claude-haiku-4-5 & 10\% & Source index, shared schemas, taxonomy skeleton, budget telemetry \\
\bottomrule
\end{tabularx}
\end{center}

\section{Wave Execution Plan}

\subsection{Wave Summary}

\begin{center}
\rowcolors{2}{rowgray}{white}
\begin{tabularx}{\textwidth}{C{1cm}L{3cm}C{2cm}C{2.5cm}X}
\toprule
\rowcolor{primary}
\tableheader{Wave} & \tableheader{Name} & \tableheader{Mode} & \tableheader{Model} & \tableheader{Outputs} \\
\midrule
0 & Foundation & Sequential & Haiku & Source index, schemas, taxonomy shell \\
1A & GL/VK + Compute & Parallel & Sonnet & Modules 01, 02 \\
1B & Amiga + XSS & Parallel & Sonnet+Opus & Modules 03, 04 \\
1C & GSD-OS + OTel & Parallel & Sonnet & Modules 05, 06 \\
2 & Synthesis & Sequential & Opus & Cross-module maps, unit circle finalization \\
3 & Publication & Sequential & Sonnet & All deliverables, compiled PDF \\
\bottomrule
\end{tabularx}
\end{center}

\subsection{Wave 0 --- Foundation}

\begin{center}
\rowcolors{2}{rowgray}{white}
\begin{tabularx}{\textwidth}{L{3cm}XC{2cm}C{2cm}}
\toprule
\rowcolor{primary}
\tableheader{Task} & \tableheader{Description} & \tableheader{Agent} & \tableheader{Tokens} \\
\midrule
Source Index & Compile all authoritative URLs, repos, and documentation references & SCOUT/Haiku & 3K \\
Shared Schemas & Define JSON schemas for example entries, translation table rows, hack catalog entries & BUDGET/Haiku & 3K \\
Taxonomy Shell & Create 1001-entry CSV skeleton with column headers; populate first 50 rows & LOG/Haiku & 10K \\
\bottomrule
\end{tabularx}
\end{center}

\subsection{Wave 1 --- Parallel Tracks}

\textbf{Track A — Modules 01 \& 02 (CRAFT-GL + CRAFT-COMPUTE / Sonnet):}

\begin{center}
\rowcolors{2}{rowgray}{white}
\begin{tabularx}{\textwidth}{L{3cm}XC{2.5cm}}
\toprule
\rowcolor{primary}
\tableheader{Task} & \tableheader{Description} & \tableheader{Tokens} \\
\midrule
Module 01 draft & Complete NeHe gap analysis; annotate 200+ examples in taxonomy & 40K \\
Module 02 draft & CUDA translation table; Pi 5 audio DSP pipeline docs & 35K \\
\bottomrule
\end{tabularx}
\end{center}

\textbf{Track B — Modules 03 \& 04 (CRAFT-AMIGA / Opus + CRAFT-XSS / Sonnet):}

\begin{center}
\rowcolors{2}{rowgray}{white}
\begin{tabularx}{\textwidth}{L{3cm}XC{2.5cm}}
\toprule
\rowcolor{primary}
\tableheader{Task} & \tableheader{Description} & \tableheader{Tokens} \\
\midrule
Module 03 draft & Full Zorro II/III spec, autoconfig, chipset DMA model, FPGA survey & 30K \\
Module 04 draft & XScreenSaver hack API, GLSL portability matrix, 50-hack catalog, Hello-Triangle skeleton & 25K \\
\bottomrule
\end{tabularx}
\end{center}

\textbf{Track C — Modules 05 \& 06 (CRAFT-GSD + CRAFT-OBS / Sonnet):}

\begin{center}
\rowcolors{2}{rowgray}{white}
\begin{tabularx}{\textwidth}{L{3cm}XC{2.5cm}}
\toprule
\rowcolor{primary}
\tableheader{Task} & \tableheader{Description} & \tableheader{Tokens} \\
\midrule
Module 05 draft & Tauri hook points, WebGL bridge, GLSL injection API, metric event bus, chipset YAML v2 & 30K \\
Module 06 draft & OTel Collector YAML, DCGM exporter config, 5 Grafana dashboard specs & 35K \\
\bottomrule
\end{tabularx}
\end{center}

\textbf{CAPCOM Gate 1:} Tibsfox reviews scope and licensing before Wave 2. Key checks: Shadertoy license compliance confirmed; all YAML configs reviewed for security; Amiga architecture claims verified against Hardware Reference Manual.

\subsection{Wave 2 --- Synthesis}

\begin{center}
\rowcolors{2}{rowgray}{white}
\begin{tabularx}{\textwidth}{L{4cm}XC{2.5cm}}
\toprule
\rowcolor{primary}
\tableheader{Task} & \tableheader{Description} & \tableheader{Tokens} \\
\midrule
Amiga$\to$GSD Mapping & Full mapping table: Zorro $\to$ DACP, autoconfig $\to$ skill autoload, chipset $\to$ YAML roles & 20K \\
Unit Circle Taxonomy & Assign $\theta$ values to all 1001 examples; produce visualization spec & 20K \\
GLSL Portability Matrix & Complete cross-version compatibility table with compat shims & 15K \\
Cross-Module Integration & Verify all six modules reference each other correctly; no orphan sections & 10K \\
\bottomrule
\end{tabularx}
\end{center}

\textbf{CAPCOM Gate 2:} Technical accuracy review. Tibsfox confirms: unit circle taxonomy assignments make sense; Amiga mapping is architecturally defensible; GSD-OS integration spec is implementable with Tauri 2.x.

\subsection{Wave 3 --- Publication and Verification}

\begin{center}
\rowcolors{2}{rowgray}{white}
\begin{tabularx}{\textwidth}{L{4cm}XC{2.5cm}}
\toprule
\rowcolor{primary}
\tableheader{Task} & \tableheader{Description} & \tableheader{Tokens} \\
\midrule
OTel YAML Validation & Lint and validate Alloy config; test against OTel Collector contrib schema & 8K \\
Hello-Triangle Skeleton & C source for a compiling XScreenSaver GLX hack with GLSL & 10K \\
Grafana Dashboard JSONs & Five complete dashboard import files & 12K \\
Bibliography Audit & Verify all URLs resolve; confirm license information; flag any 404s & 8K \\
Final LaTeX Compilation & Three-pass xelatex; PDF + .tex source & 5K \\
Index Page & index.html with download cards & 5K \\
\bottomrule
\end{tabularx}
\end{center}

\subsection{Cache Optimization Strategy}

\begin{itemize}[noitemsep]
  \item Wave 0 shared schemas are computed once and injected into all six Wave 1 agent contexts as a shared prefix — exploiting the 5-minute cache TTL
  \item The source index (URLs, repos, documentation references) is injected into all agents, eliminating redundant web research
  \item The 1001-example taxonomy shell CSV is computed in Wave 0 and appended to by each Wave 1 track, avoiding the $O(n^2)$ re-read problem
  \item Modules 03 and 04 share the Amiga Hardware Reference Manual excerpt as a common context block
\end{itemize}

\subsection{Token Budget}

\begin{center}
\rowcolors{2}{rowgray}{white}
\begin{tabularx}{\textwidth}{L{3cm}C{2.5cm}C{2.5cm}C{2.5cm}}
\toprule
\rowcolor{primary}
\tableheader{Wave} & \tableheader{Opus Tokens} & \tableheader{Sonnet Tokens} & \tableheader{Haiku Tokens} \\
\midrule
Wave 0 & 0 & 0 & 23K \\
Wave 1A & 0 & 75K & 0 \\
Wave 1B & 30K & 25K & 0 \\
Wave 1C & 0 & 65K & 0 \\
Wave 2 & 65K & 0 & 0 \\
Wave 3 & 0 & 48K & 0 \\
\textbf{Total} & \textbf{95K} & \textbf{213K} & \textbf{23K} \\
\textbf{Split} & \textbf{29\%} & \textbf{64\%} & \textbf{7\%} \\
\bottomrule
\end{tabularx}
\end{center}

\subsection{Dependency Graph}

\begin{asciibox}
Source Index ──────────────────┐
Shared Schemas ─────────────┐  │
Taxonomy Shell ──────────┐  │  │
                         │  │  │
         ┌───────────────┘  │  │
         ▼                  │  │
    Module 01 (GL/VK) ──────┤  │
    Module 02 (Compute) ────┤  │
    Module 03 (Amiga) ──────┤  │  ← Opus priority
    Module 04 (XSS) ────────┤  │
    Module 05 (GSD-OS) ─────┤  │
    Module 06 (OTel) ───────┘  │
         │                     │
         ▼                     │
    Wave 2 Synthesis ──────────┘
         │
         ▼
    Wave 3 Publication
         │
         ▼
    All Deliverables (D1─D9)
\end{asciibox}

\section{Test Plan}

\subsection{Test Categories}

\begin{center}
\rowcolors{2}{rowgray}{white}
\begin{tabularx}{\textwidth}{L{3cm}XC{2cm}C{2.5cm}C{1.5cm}}
\toprule
\rowcolor{primary}
\tableheader{Category} & \tableheader{Purpose} & \tableheader{Count} & \tableheader{Failure Action} & \tableheader{Priority} \\
\midrule
Safety-Critical & Non-negotiable boundaries & 6 & BLOCK & Mandatory \\
Core Functionality & Deliverable acceptance & 22 & BLOCK & Required \\
Integration & Cross-module validation & 10 & BLOCK & Required \\
Edge Cases & Robustness & 7 & LOG & Best-effort \\
\textbf{Total} & & \textbf{45} & & \\
\bottomrule
\end{tabularx}
\end{center}

\subsection{Safety-Critical Tests}

\begin{center}
\rowcolors{2}{rowgray}{white}
\begin{tabularx}{\textwidth}{C{1.5cm}L{4cm}X}
\toprule
\rowcolor{primary}
\tableheader{ID} & \tableheader{Verifies} & \tableheader{Expected} \\
\midrule
SC-SRC & Source quality & All citations traceable to Khronos, Raspberry Pi Foundation, Grafana Labs, or equivalent professional sources. Zero forum posts as primary sources. \\
SC-NUM & Numerical attribution & Every GFLOPS figure, latency measurement, and percentage attributed to a specific benchmark, datasheet, or publication. \\
SC-LIC & Open-source licensing & All Shadertoy-derived examples verified CC0/MIT/BSD/CC-BY only. No CC-BY-NC-SA content in GSD-OS distribution. \\
SC-SEC & Security: no hardcoded credentials & OTel config YAML uses environment variable references for all API tokens. Zero literal secrets in any deliverable. \\
SC-HW & Hardware safety: QPU DMA warning & VC4CL section explicitly documents root privilege requirement and memory management unit absence on VideoCore IV. \\
SC-ADV & No policy advocacy & Document presents GPU architecture comparisons without vendor advocacy. CUDA/OpenCL comparison presented as technical trade-off analysis. \\
\bottomrule
\end{tabularx}
\end{center}

\subsection{Core Functionality Tests (Selected)}

\begin{center}
\rowcolors{2}{rowgray}{white}
\begin{tabularx}{\textwidth}{C{1.5cm}XL{5cm}}
\toprule
\rowcolor{primary}
\tableheader{ID} & \tableheader{Verifies} & \tableheader{Expected} \\
\midrule
CF-01 & Example taxonomy completeness & All 1001 entries have: source URL, API level, technique tag, Pi-5 compat, GSD applicability \\
CF-02 & NeHe gap bridged & All 48 NeHe lessons have a modern equivalent with GLSL/SPIR-V source \\
CF-03 & CUDA translation table coverage & $\geq$ 7 CUDA patterns with OpenCL/Vulkan/NEON equivalents \\
CF-04 & Pi 5 audio pipeline documented & End-to-end path from ALSA/JACK to NEON/OpenCL kernel to output with latency spec \\
CF-05 & Amiga autoconfig documented & Zorro II/III autoconfig mechanism fully specified with GSD DACP mapping \\
CF-06 & Chipset mapping complete & Agnus/Denise/Paula each mapped to GSD chipset YAML role with rationale \\
CF-07 & Hello-Triangle hack compiles & XScreenSaver skeleton builds with \texttt{gcc} and runs as GLX screensaver \\
CF-08 & 50-hack catalog complete & All 50 entries have: name, source, GLSL version, GPU load estimate, license \\
CF-09 & OTel YAML validates & Alloy config lints without errors against OTel Collector contrib schema \\
CF-10 & All 5 dashboards specified & GPU Health, Pi SoC, Audio DSP Latency, GSD Skill Metrics, Pipeline Health \\
\bottomrule
\end{tabularx}
\end{center}

\subsection{Verification Matrix}

\begin{center}
\rowcolors{2}{rowgray}{white}
\begin{tabularx}{\textwidth}{XL{3.5cm}C{1.5cm}}
\toprule
\rowcolor{primary}
\tableheader{Success Criterion} & \tableheader{Test ID(s)} & \tableheader{Status} \\
\midrule
1001 examples tagged and cross-referenced & CF-01, CF-02 & Pending \\
NeHe gap bridged with modern equivalents & CF-02 & Pending \\
CUDA$\to$OpenCL translation table & CF-03 & Pending \\
Pi 5 audio DSP pipeline end-to-end & CF-04 & Pending \\
Amiga autoconfig $\to$ DACP mapping & CF-05 & Pending \\
Chipset coprocessor model $\to$ GSD YAML & CF-06 & Pending \\
XScreenSaver Hello-Triangle skeleton & CF-07 & Pending \\
50-hack curated GSD-OS catalog & CF-08 & Pending \\
GSD-OS Tauri integration specified & IN-01, IN-02 & Pending \\
OTel pipeline configured end-to-end & CF-09, SC-SEC & Pending \\
5 Grafana dashboards specified & CF-10 & Pending \\
All cross-module references consistent & IN-03, IN-04 & Pending \\
\bottomrule
\end{tabularx}
\end{center}

\section{Mission Crew Manifest}

\begin{center}
\rowcolors{2}{rowgray}{white}
\begin{tabularx}{\textwidth}{L{2cm}L{3.5cm}XC{2cm}}
\toprule
\rowcolor{primary}
\tableheader{Role} & \tableheader{Agent} & \tableheader{Responsibility} & \tableheader{Model} \\
\midrule
FLIGHT & Orchestrator & Mission direction; wave go/no-go gates; budget approval & Opus \\
PLAN & Planner & Wave decomposition; dependency graph; parallel track management & Opus \\
TOPO & Topology Manager & Mid-mission agent restructuring; team-of-teams coordination & Opus \\
ARCH & Architecture & Cross-cutting decisions: taxonomy schema, YAML config structure & Opus \\
ANALYST & Pattern Analyst & Unit circle taxonomy $\theta$ assignments; cross-module pattern detection & Opus \\
INTEG & Integration & Cross-module reference validation; Amiga$\to$GSD mapping synthesis & Opus \\
RETRO & Retrospective & Post-wave analysis; lessons learned for skill-creator & Opus \\
CRAFT-GL & OpenGL/Vulkan Specialist & Modules 01 \& 02 primary author & Sonnet \\
CRAFT-AMIGA & Amiga Architecture Specialist & Module 03 primary author; Zorro/DMA domain authority & Sonnet \\
CRAFT-XSS & XScreenSaver/GLSL Specialist & Module 04; Hello-Triangle skeleton; GLSL portability & Sonnet \\
CRAFT-GSD & GSD-OS Integration Specialist & Module 05; Tauri hook points; WebGL bridge & Sonnet \\
CRAFT-OBS & Observability Specialist & Module 06; OTel YAML; Grafana dashboards & Sonnet \\
EXEC-1 & Executor Track A & Document production: Modules 01, 02 & Sonnet \\
EXEC-2 & Executor Track B & Document production: Modules 03, 04 & Sonnet \\
EXEC-3 & Executor Track C & Document production: Modules 05, 06 & Sonnet \\
SECURE & Security & License compliance; credential scanning; hardware safety notes & Sonnet \\
VERIFY & Verification & Source validation; numerical attribution; URL audit & Sonnet \\
CAPCOM & HITL Interface & Only channel to Tibsfox; gate approvals; scope change requests & Sonnet \\
SCOUT & Research Agent & Source gathering; web research; filling identified gaps & Haiku \\
BUDGET & Resource Manager & Token budget tracking; session health telemetry & Haiku \\
LOG & Chronicle & Commit messages; wave summaries; audit trail & Haiku \\
\bottomrule
\end{tabularx}
\end{center}

\section{Execution Summary}

\begin{center}
\rowcolors{2}{rowgray}{white}
\begin{tabularx}{\textwidth}{L{5cm}X}
\toprule
\rowcolor{primary}
\tableheader{Metric} & \tableheader{Value} \\
\midrule
Activation Profile & Fleet (21 roles) \\
Research Modules & 6 \\
Wave Count & 4 (Wave 0 + Wave 1 parallel ×3 + Wave 2 + Wave 3) \\
Parallel Tracks & 3 simultaneous in Wave 1 \\
Estimated Context Windows & 18--24 across 6--8 sessions \\
Total Token Budget (est.) & 331K (Opus: 95K / Sonnet: 213K / Haiku: 23K) \\
Model Split & 29\% Opus / 64\% Sonnet / 7\% Haiku \\
CAPCOM Gates & 2 (post Wave 1, post Wave 2) \\
Deliverables & 9 (D1--D9) \\
Total Tests & 45 (6 safety-critical / 22 core / 10 integration / 7 edge) \\
Success Criteria & 12 \\
Primary Authorities & Khronos Group, Raspberry Pi Foundation, Grafana Labs, jwz.org \\
\bottomrule
\end{tabularx}
\end{center}

\vspace{2em}

\begin{quotebox}
\textit{``The Amiga did not win by being the most powerful machine of 1985. It won by being the most \textbf{architecturally honest} machine of 1985 — each coprocessor doing exactly the one thing it was built to do, exactly when it was supposed to do it, without interrupting the others. The spaces between their DMA cycles were where the music played.''}

\vspace{0.8em}
\noindent This mission follows that architecture from the fragment shader to the Grafana dashboard. Every example in the 1001-example taxonomy is a clock cycle. The understanding lives in the spaces between them. Hand this document to GSD skill-creator and let the coprocessors run.
\end{quotebox}

\end{document}
